%% SCRAP: architecture/03-architecture/hal/README
%% SOURCE: docs/working/architecture/03-architecture/hal/README.md
%% STATUS: WORKING
%% FITS: dev-guide/ch-platform
%% EDITORIAL: lifted — prose rewritten to press voice

\section{The Hardware Abstraction Layer: Orientation}

The Hardware Abstraction Layer (HAL) is the architectural seam that lets a
single StarForth source base run as a hosted virtual machine, as a microkernel
guest, and as the bare-metal StarKernel without forking the interpreter. This
section orients the reader to the HAL document set and the roles it serves.

\subsection{Reading Order}

The HAL documentation is layered. Each document answers a distinct question and
builds on the one before it.

\begin{itemize}
  \item \textbf{Overview} --- the problem the HAL solves and how it positions
        StarForth within the StarForth $\rightarrow$ StarKernel $\rightarrow$
        StarshipOS progression. It covers the three-layer model, the design
        principles (VM purity, contract-first interfaces, testability), the HAL
        subsystems, and the relationship between the HAL and the physics
        subsystems.
  \item \textbf{Interfaces} --- the contract. Precise function signatures,
        guaranteed semantics, error handling, performance envelopes, and the
        concurrency model (ISR-safe versus thread-safe) for each subsystem:
        time and timers, interrupts, memory, console, CPU, and panic.
  \item \textbf{Platform Implementations} --- how to satisfy the contract on a
        concrete target, with complete worked examples for hosted Linux and
        freestanding StarKernel, plus a platform testing strategy and common
        pitfalls.
  \item \textbf{Migration Plan} --- the staged refactor that moves the existing
        code base onto the HAL with zero functional regressions.
  \item \textbf{StarKernel Integration} --- kernel-specific detail: the UEFI
        boot sequence, the freestanding C environment, hardware bring-up, and
        the path to a working \texttt{ok} prompt.
\end{itemize}

\subsection{Roles}

The document set serves four audiences. A VM developer reads the overview and
the interfaces, then calls HAL functions and never touches a platform API
directly. A platform author adds a new target reads through the implementation
guide and satisfies every interface, validating against the full VM test suite.
A core developer performing the migration follows the staged plan and tests
after each step. A kernel developer reads everything, with emphasis on the
integration guide, and builds incrementally from UEFI boot to the REPL prompt.

\subsection{Governing Principles}

The HAL rests on five principles. \emph{VM purity}: interpreter code is
platform-agnostic and never learns which platform it runs on. \emph{Contract-first}:
interfaces are contracts with precise semantics, not convenience wrappers.
\emph{Testability}: kernel-bound code is developed and tested on Linux before it
ships to bare metal. \emph{Zero overhead}: optimized builds inline HAL calls to
direct hardware access. \emph{Fail-fast}: platform assumptions are validated at
initialization, not during execution.

\subsection{Success Criteria}

The HAL is judged successful when the VM core carries zero platform-specific
code, all of the more than nine hundred regression tests pass on every
platform, the runtime's zero algorithmic variance is preserved across
platforms, no measurable performance regression is introduced, StarKernel boots
to its \texttt{ok} prompt, and the physics subsystems behave identically
everywhere.

%% TODO(bob): confirm exact regression-test count for the published edition
%%   (source README cites "936+"; CLAUDE.md cites "800/800" for ACL POST and
%%   "936+" for the full suite).
